\section{Related Work}

\subsection{Ad-Hoc Data Acquisition and Organization in Physical Computer Vision}
The rapid proliferation of deep learning across scientific domains has driven the creation of numerous domain-specific image datasets. However, empirical dataset collection practices in applied computer vision have historically prioritized sample volume over methodological data governance~\cite{varoquaux2022}. In typical dataset generation pipelines, physical samples are photographed, cropped into sub-images or patches, and stored in unindexed directory trees organized solely by target class labels. Important provenance metadata---such as physical sample identifiers, harvesting locations, acquisition timestamps, camera sensor settings, and surface preparation conditions---are rarely recorded or embedded into machine-readable schemas.

This ad-hoc approach creates severe vulnerabilities during dataset partitioning. Because standard machine learning frameworks (e.g., PyTorch `ImageFolder` or Scikit-Learn `train_test_split`) operate strictly on file paths, dataset splits are constructed via uniform random sampling over individual image files. When multiple image observations originate from the same physical entity, this unindexed organization guarantees that observations from the same physical source are distributed across training and evaluation splits, corrupting the scientific integrity of model evaluation~\cite{kapoor2023}.

\subsection{Data Leakage, Deployment Gap, and Same-Specimen-Picture Bias (SSPB)}
Data leakage is increasingly recognized as a primary driver of the reproducibility crisis across computational sciences~\cite{kapoor2023,kaufman2012}. Kaufman et al.~\cite{kaufman2012} established an early taxonomy of leakage in data mining, distinguishing between feature leakage and evaluation leakage. Kapoor and Narayanan~\cite{kapoor2023} conducted a systematic review of 329 machine learning papers across 17 scientific disciplines, discovering that over 50\% exhibited data leakage, leading to widespread over-claiming of predictive performance.

In botanical image analysis and wood species recognition, **Same-Specimen-Picture Bias (SSPB)** was explicitly conceptualized by Figueroa-Mata et al.~\cite{figueroamata2022} in their identification of Costa Rican tree species, emphasizing that multiple image captures from the same physical wood specimen introduce severe evaluation bias if not isolated during data splitting. Rosa da Silva et al.~\cite{rosadasilva2022} further demonstrated that evaluation protocols routinely ignore multiple image captures per tree across anatomical planes, yielding overoptimistic performance estimates.

In timber forensics, Ravindran et al.~\cite{ravindran2019,ravindran2020,ravindran2021,ravindran2022} and Wiedenhoeft et al.~\cite{wiedenhoeft2011} documented a substantial "deployment gap"---reporting accuracy drops ranging from approximately 10\% (e.g., Ravindran et al.~\cite{ravindran2019} evaluating Ghanaian timber models on field samples versus laboratory xylarium specimens) up to 25.0\% under real-world field conditions. 

Crucially, our work differs from prior deployment gap studies in a fundamental methodological aspect: while prior field deployment studies conflated specimen leakage with external camera domain shift, sensor variation, and field environmental noise, our framework isolates **pure specimen-level data leakage** under strictly controlled experimental conditions (matched feature backbones, identical classifiers, and identical image acquisition distributions).

In medical computer vision, Roberts et al.~\cite{roberts2021} reviewed over 400 deep learning models developed for COVID-19 detection from chest radiographs, concluding that none were suitable for clinical deployment due to patient-level data leakage. Varoquaux and Cheplygina~\cite{varoquaux2022} further demonstrated that machine learning benchmarks in medical imaging routinely overestimate real-world diagnostic utility when group independence is violated. In explainable AI (XAI), Geirhos et al.~\cite{geirhos2020} and Lapuschkin et al.~\cite{lapuschkin2019} showed that deep neural networks under leaked splits act as "Clever Hans" predictors, exploiting non-taxonomic background shortcuts rather than learning true domain features.

\subsection{Group-Aware Data Partitioning Protocols}
To prevent group leakage, specialized partitioning strategies have been explored in specific subfields. Standard machine learning libraries provide basic utilities such as `GroupKFold`~\cite{scikit}, which enforces group isolation across splits. However, naive GroupKFold treats groups homogeneously, ignoring per-class sample imbalance and representation density, which often results in missing minority classes in test splits.

In spatial ecology, Roberts et al.~\cite{roberts2017} introduced spatial and phylogenetic blocking to mitigate spatial autocorrelation in species distribution modeling. In metallography, Azimi et al.~\cite{azimi2018} demonstrated the necessity of specimen-level splitting for alloy microstructure classification. In fine-grained wood identification, recent studies on taxonomically close species (e.g., Liu et al.~\cite{liu2025} for *Pterocarpus* species identification) and regional timber conservation tools (e.g., Song et al.~\cite{song2025} for Vietnamese timber identification) have increasingly adopted specimen-disjoint cross-validation protocols to ensure realistic evaluation on unseen timber blocks.

\subsection{Automated Wood Identification and Timber Forensics}
Macroscopic wood species identification plays a pivotal role in enforcing CITES Appendix II regulations and combating illegal timber trafficking~\cite{cites,dormontt2015}. Recent computer vision research has transitioned from handcrafted texture descriptors (GLCM, LBP, Gabor filters)~\cite{woodreview} to deep convolutional networks and vision transformers~\cite{fabijanska2021,wu2021,ravindran2022,song2025}.

Despite architectural progress, wood species identification presents extreme fine-grained challenges~\cite{wiedenhoeft2011}. Taxonomically close sibling species within the same genus (e.g., \textit{Afzelia}, \textit{Dalbergia}, or \textit{Pterocarpus}~\cite{liu2025}) share near-identical macroscopic anatomical traits, including vessel pore distributions and axial parenchyma bands. Evaluating models on random image splits obscures these biological challenges by allowing models to memorize specimen-specific surface shortcuts. A standardized data governance framework is therefore essential to establish reliable, leak-free benchmarks for forensic timber identification.
